\section{What the residue section selects in the unchanged
source}\label{what-the-residue-section-selects-in-the-unchanged-source}

25 September 2026. Complete derivation RSS0--RSS6. This calculation
tests an agent-selected observation in the established coefficient
source. It changes none of the user's definitions. The supporting datum
remains \(\tau\langle Z_1;\mathrm{no}\ Z_2\rangle\); its addition
remains retracted. All scalar operations occur after complete arithmetic
reconstruction, with branch histories retained.

\subsection{RSS0. The source and the precise
comparison}\label{rss0.-the-source-and-the-precise-comparison}

Keep the entire original source and its topology: \[
\mathcal B=\{F\in\mathcal O(\mathbb C):q_{A,N}(F)<\infty\},\qquad
q_{A,N}(F)=\sup_{|\Re s|\le A}(1+|\Im s|)^N|F(s)|,
\tag{RSS0.1}
\] for all \(A>0\) and integers \(N\ge0\). The closed ideal \(I\)
requires every vanishing jet through order \(m_\rho-1\) at every
original nontrivial zero \(\rho\) of \(\zeta\). Retain \[
F_0(s)=\frac{s(s-1)}8\pi^{-s/2}\Gamma(s/2)\zeta(s),\quad
F_0(0)=\frac18,\quad g_t(s)=e^{ts^2},\quad t>0,
\tag{RSS0.2}
\] with the original exceptional-point comparisons proved in NPE0 and
the earlier full source. No factor is being discarded or substituted for
original zeta. Write \[
h_t=8g_tF_0/s,\qquad j_t=g_t/s,\qquad
E=\mathcal B+\mathbb C h_t=\mathcal B+\mathbb C j_t.
\tag{RSS0.3}
\] NPE1 and NER5 prove these are the same actual meromorphic functions
with the same topology. Their difference \[
k_t=j_t-h_t=\frac{g_t(1-8F_0)}s\in\mathcal B
\tag{RSS0.4}
\] is retained. The numerator vanishes at zero; its quotient has rapid
strip decay. The action throughout is \(U_n f(s)=n^{1-s}f(s)\), for all
recovered positive integers.

The earlier NCI theorem and NPE10 say that choosing the
zero-\(h_t\)-coordinate dual slice recovers \(I^\perp\) as its largest
all-cover-invariant subspace. We now calculate the selection made by an
entire family of different residue sections in this same \(E\). These
are comparisons with the original observation, not replacements for it.

\subsection{RSS1. The Gaussian-section discrepancies span the entire
source}\label{rss1.-the-gaussian-section-discrepancies-span-the-entire-source}

Define \[
\eta_{n,t}=nj_t-U_nj_t
=g_t(s)\frac{n-n^{1-s}}s\in\mathcal B.
\tag{RSS1.1}
\] The divided value at zero is exactly \(n\log n\). We prove \[
\boxed{\overline{\operatorname{span}\{\eta_{n,t}:n\ge2\}}^{\mathcal B}=\mathcal B.}
\tag{RSS1.2}
\] Let \(\Lambda\in\mathcal B'\) annihilate the displayed family. The
source-valued entire function \[
K_w(s)=g_t(s)\frac{e^w-e^{(1-s)w}}s
=we^w g_t(s)\int_0^1 e^{-\theta sw}\,d\theta
\tag{RSS1.3}
\] has removable value \(we^w\) at zero. On \(|\Re s|\le A\), writing
\(w=a+ib\), the integral expression is bounded using \[
e^{-ty^2+\theta by}\le
e^{b^2/(4t)}e^{-t(y-\theta b/(2t))^2},\qquad0\le\theta\le1.
\tag{RSS1.4}
\] Every vertical polynomial is controlled by a polynomial in \(|b|\)
times a Gaussian maximum. Differentiation on compact \(w\)-sets
preserves these bounds, proving source holomorphy. Continuity of
\(\Lambda\) bounds it by finitely many source seminorms, so the entire
scalar function \(L(w)=\Lambda(K_w)\) satisfies
\(|L(w)|\le C e^{C(1+|w|^2)}\). It vanishes at all \(\log n\), including
\(0\). A nonzero entire function with this bound has at most \(O(R^2)\)
zeros in a radius-\(R\) disk centered at a fixed point \(w_0\) with
\(L(w_0)\ne0\), by Jensen's formula with outer radius \(2R\). The
distinct \(\log n\) for \(n\le e^{R/2}\) violate that bound. Thus
\(L=0\).

The exact derivative is \[
(\partial_w-1)K_w=g_te^{(1-s)w}.
\tag{RSS1.5}
\] Consequently \(\Lambda(g_te^{vs})=0\) for every complex \(v\). For
completeness the passage from these tests to the entire source retains
the constants. Given \(F\in\mathcal B\) and real \(u>t\), put
\(H=e^{(u-t)s^2}F\). Its inverse and forward half-Mellin expressions are
\[
a_H(x)=\frac{x^{-1/2}}\pi\int_{\mathbb R}H(1/2+iy)x^{-iy}\,dy,
\qquad H(s)=\frac12\int_{\mathbb R}a_H(e^v)e^{sv}\,dv.
\tag{RSS1.6}
\] These are the original Fourier inversion with the factor \(1/2\)
retained. Shifting the integration line to arbitrary fixed real parts
using rapid strip decay proves that \(a_H(e^v)\) decreases faster than
every exponential in \(|v|\). The second integral, after multiplication
by \(g_t\), therefore converges in every \(\mathcal B\) seminorm:
\(q_{A,N}(g_te^{vs})\le C_{A,N,t}e^{A|v|}\) for real \(v\). Applying
\(\Lambda\) gives \(\Lambda(e^{us^2}F)=0\) for every real \(u>t\). The
function of \(u\) is holomorphic on \(\Re u>0\), by the uniform negative
vertical quadratic on compact parameter sets. The identity theorem
extends its vanishing throughout that half-plane. Finally the exact
integral remainder gives \[
q_{A,N}((e^{us^2}-1)F)
\le u e^{A^2}(A^2+1)q_{A,N+2}(F),\qquad0<u\le1.
\tag{RSS1.7}
\] Letting \(u\downarrow0\) proves \(\Lambda(F)=0\). Hahn--Banach
separation of closed linear subspaces now proves (RSS1.2). No
finite-height zero computation or RH assumption enters.

\subsection{RSS2. Arbitrary polynomial sections and their complete
vanishing
ideals}\label{rss2.-arbitrary-polynomial-sections-and-their-complete-vanishing-ideals}

Take any complex polynomial \(P\) with \(P(0)=1\), including the
constant polynomial \(1\), and define \[
h_{P,t}=g_tP/s,\qquad
h_{P,t}-j_t=g_t(P-1)/s\in\mathcal B.
\tag{RSS2.1}
\] Thus every such section belongs to the same actual \(E\), has residue
one, and gives the same product topology. It does not change the source
or action. Its discrepancy is exactly \[
\delta^P_{n,t}=nh_{P,t}-U_nh_{P,t}=P\eta_{n,t}.
\tag{RSS2.2}
\] Let \(a\) run over the distinct zeros of \(P\), with orders \(d_a\).
None is zero. Polynomial multiplication is continuous on \(\mathcal B\):
each strip seminorm is controlled by one with its polynomial weight
increased by \(\deg P\). Its image satisfies \[
P\mathcal B=\{F\in\mathcal B:F^{(j)}(a)=0\text{ for }0\le j<d_a\}.
\tag{RSS2.3}
\] One inclusion follows from the product rule. For the other, these
vanishing orders make \(F/P\) entire. Choose small disjoint disks about
the finitely many roots. Outside their union, \(1/P\) is bounded: this
follows from its polynomial behavior at infinity and its nonzero minimum
on the remaining compact region. Inside each disk, the maximum principle
bounds \(F/P\) by its boundary values. Those boundaries lie in a fixed
compact set where \(1/P\) is bounded. Thus for each \(A,N\), \[
q_{A,N}(F/P)\le C_P q_{A,N}(F)+C_{P,A,N}q_{A_0,0}(F)
\tag{RSS2.4}
\] for a sufficiently large fixed \(A_0\). This proves membership and
continuity of division on the image. The right side of (RSS2.3) is an
intersection of kernels of finitely many continuous jet evaluations,
hence is closed.

By (RSS1.2), every \(F\in\mathcal B\) is a limit of linear combinations
of the \(\eta_{n,t}\). Multiplication by \(P\) preserves that
convergence. The closed-image identity therefore proves \[
\boxed{\overline{\operatorname{span}\{\delta^P_{n,t}:n\ge2\}}^{\mathcal B}=P\mathcal B.}
\tag{RSS2.5}
\] This is a global theorem for every displayed polynomial, rather than
a bounded collection of examples.

\subsection{RSS3. Exactly what the zero-coordinate observation
selects}\label{rss3.-exactly-what-the-zero-coordinate-observation-selects}

In the \(h_{P,t}\) presentation write \[
\widetilde\Lambda(F+c h_{P,t})=\Lambda(F)+c\alpha_P.
\] The exact transpose action is \[
(U_n^E)'(\Lambda,\alpha_P)
=(U_n'\Lambda,n\alpha_P-\Lambda(\delta^P_{n,t})).
\tag{RSS3.1}
\] Therefore every invariant linear subspace contained in the slice
\(\alpha_P=0\) is contained in \((P\mathcal B)^\perp\oplus\{0\}\), by
(RSS2.5). Conversely \(P\mathcal B\) is invariant under every \(U_n\),
since multiplication commutes, so its annihilator with zero coordinate
is invariant. This proves the exact maximality statement \[
\boxed{\text{largest all-cover-invariant subspace in }\alpha_P=0
=(P\mathcal B)^\perp\oplus\{0\}.}
\tag{RSS3.2}
\] For \(P=1\) this is \(\{0\}\). For a general \(P\), its first factor
is exactly the span of the full jet functionals
\(\operatorname{ev}_a^{(j)}\), \(0\le j<d_a\). To prove that assertion,
the finite jet map from \(\mathcal B\) onto
\(\bigoplus_a\mathbb C^{d_a}\) is surjective: multiply a polynomial
interpolation function by \(e^{s^2}\). The latter factor is a local
unit, so its triangular jet action is invertible at every prescribed
point; finite polynomial interpolation realizes the inverse prescribed
jets. One proof of that interpolation is the Chinese remainder theorem
for the pairwise coprime ideals \((s-a)^{d_a}\) in \(\mathbb C[s]\),
obtained by Bezout identities and their resulting idempotents. The
kernel is (RSS2.3). Hence every annihilating functional factors through
precisely those finite jets.

The actual action on these jets retains every lower term: \[
U_n'\operatorname{ev}_a^{(j)}
=n^{1-a}\sum_{b=0}^j\binom jb(-\log n)^{j-b}\operatorname{ev}_a^{(b)}.
\tag{RSS3.3}
\] This follows by differentiating \(n^{1-s}F(s)\). It works at every
nonzero complex \(a\) allowed by the polynomial section. Thus this
observation and cover invariance alone do not select a critical-line
divisor. The equation is an explicit comparison with the original
source, not an assertion that these alternative sections are compatible
with all of the programme's geometric data.

\subsection{RSS4. The original zeta sector remains an exact
graph}\label{rss4.-the-original-zeta-sector-remains-an-exact-graph}

For the Gaussian section put \(k=k_t\). For the same actual functional,
its old and new coordinates satisfy \[
\beta=\widetilde\Lambda(j_t)=\alpha+\Lambda(k).
\tag{RSS4.1}
\] This continuous invertible shear carries the original zero-coordinate
lift to \[
\{(\Lambda,\Lambda(k))_j:\Lambda\in I^\perp\}.
\tag{RSS4.2}
\] It is invariant. Indeed \(\eta_n=\delta_n+(n-U_n)k\), where the
original \(\delta_n=8F_0\eta_n\) lies in \(I\). For
\(\Lambda\in I^\perp\), \[
n\Lambda(k)-\Lambda(\eta_n)=\Lambda(U_nk)=(U_n'\Lambda)(k).
\tag{RSS4.3}
\] This is exactly the last coordinate required by (RSS4.2). The old
sector has not vanished when the new zero-coordinate slice is zero. They
are different subspaces, related by the proved shear.

The full receiver retains the same observation as well. With
\(\phi_0=-1/s\) and \(\phi_m=(m^{1-s}-(m+1)^{1-s})/s\) for \(m\ge1\),
the corrected entire tests in the two presentations are
\(g_t\phi_m+h_t\) and \(g_t\phi_m+j_t\). Consequently \[
H_j\Lambda=H_h\Lambda+\frac{\overline{\Lambda(k)}}{1-z},\qquad
M_t\widetilde\Lambda
=H_h\Lambda-\frac{\overline\alpha}{1-z}
=H_j\Lambda-\frac{\overline\beta}{1-z}.
\tag{RSS4.4}
\] Each series converges in the compact-open topology of the disk, by
the original source seminorm estimates. The same formula carries the
full Hardy graph and its norm exactly. No conclusion about the
individual summands' Hardy membership follows by dropping their
difference.

\subsection{RSS5. The original quotient and all
multiplicities}\label{rss5.-the-original-quotient-and-all-multiplicities}

The preceding section family is compared with the actual zeta result,
not substituted for it: \[
\overline{\operatorname{span}\{8F_0\eta_{n,t}:n\ge2\}}^{\mathcal B}
=\overline{F_0\mathcal B}^{\mathcal B}=I,
\qquad F_0\mathcal B\subsetneq I.
\tag{RSS5.1}
\] NCI2--5 proves the nonpolynomial division and closure with the full
Hadamard product, clustered-disk bounds and every original multiplicity;
NER11 checks that proof. These inputs are necessary here. They are not
replaced by the finite-polynomial argument in RSS2. The exact original
section is thereby what inserts this particular divisor into the
zero-coordinate observation. The global density calculation proves
recovery of that full divisor, not a bound on its location.

For any finite subcollection of the actual nontrivial zeros with
selected orders \(1\le d_a\le m_a\), the polynomial \[
P(s)=\prod_a(1-s/a)^{d_a}
\tag{RSS5.2}
\] has \(P(0)=1\) and \(I\subset P\mathcal B\). The identity on
\(\mathcal B\) therefore induces the exact continuous surjection \[
\mathcal B/I\longrightarrow\mathcal B/(P\mathcal B),\qquad
[F]_I\longmapsto[F]_{P\mathcal B}.
\tag{RSS5.3}
\] Its kernel is \(P\mathcal B/I\). Surjectivity follows from
representatives; continuity and the quotient topology follow from the
composition of the two quotient maps. On the dual, it is precisely
inclusion of the selected full jets into \(I^\perp\), with (RSS3.3)
unchanged. Thus the family has an exact receiving map from the original
divisor, including multiplicity and topology.

\subsection{RSS6. The foundational conclusion actually
established}\label{rss6.-the-foundational-conclusion-actually-established}

The underlying meromorphic source, its cover action, residue sequence
and full receiver were kept fixed. Varying only the chosen residue
section changes the maximal invariant part of its zero-coordinate
observation from the whole original divisor dual to zero or to any
finite prescribed nonzero-point jet dual. Their source maps and all
changes are explicit above. It follows that the zero-coordinate
restriction is substantive mathematical data. Treating it as an
intrinsic property of the whole source would be an agent-introduced
error.

The original zeta section and full ideal remain available, and their
shear into the alternative presentation preserves every original
functional. A further geometric argument must operate on that retained
source and its actual specialization map. Its purity cannot be supplied
by choosing a section whose zeros were prescribed. This calculation
identifies precisely which part of the current coefficient recovery is
supplied by the chosen correction, while proving the complete morphisms
needed to return to the original construction.

Sources actually used: complete NPE0--10, NER0--11, NCI0--8 and SSR0--7
programme proofs. Their full texts accompany this continuation. The
human source of the Hardy coefficient family remains S. Waleed Noor,
\href{https://arxiv.org/abs/1809.09577v4}{arXiv:1809.09577v4}; no new
whole-paper reading or novelty claim for the underlying classical
interpolation is made. The geometric comparison remains Deligne
\href{https://numdam.org/item/PMIHES_1980__52__137_0/}{Weil II §3.6} and
Connes--Consani \href{https://arxiv.org/abs/0903.2024v3}{§5}. No
original lifting vanishing or RH theorem is asserted.
